\chapter{预期达到的目标}
\section{预期完成的任务}
本课题的主要目标是：1.证明联邦-分割学习在非独立同分布数据和标签噪声存在时的收敛性。2.设计同时对于非独立同分布数据和标签噪声具有鲁棒性的联邦-分割学习算法。

首先，通过分析联邦-分割学习在非独立同分布数据环境中的行为，并对标签噪声对数据质量和分布的影响进行建模，完成联邦-分割学习在非独立同分布数据和标签噪声存在时的收敛性证明。该收敛性证明将有助于理解和描述联邦-分割学习在不理想数据条件下的训练过程，为后续的算法改进提供理论支持。

其次，本文将设计一种同时针对非独立同分布数据和标签噪声具有鲁棒性的联邦-分割学习算法。为此，本文计划引入新的损失函数和聚合策略，旨在提升模型在面对不一致标签时的表现。此外，结合鲁棒性学习理论和自蒸馏技术，对模型的更新过程进行优化，使其更好地适应具有标签噪声的复杂数据环境。

本研究的创新点主要体现在以下几个方面。首先，在联邦-分割学习上进行系统的收敛性分析，填补当前相关研究的空白。其次，本文将利用主流的数据建模框架，量化评估标签噪声对数据质量和数据分布的影响，从而为收敛性分析和联邦-分割学习算法改进提供充分的依据。最后，本文将设计出一种能够同时应对非独立同分布数据和噪声标签的鲁棒性算法，该新算法将为分布式学习领域提供一种在实际应用中对全新解决方案。

\section{预期取得的研究成果}
预期在高水平学术会议或期刊上发表2-3篇论文。
目前，我作为第三作者的一篇论文已被NeurIPS 2024接收，我作为第一作者的一篇论文正在INFOCOM 2025审稿中。